\documentclass[10pt,a4paper]{amsart}
\usepackage[margin=19mm]{geometry}
\usepackage{amsmath,amssymb,mathtools}
\usepackage[scr=rsfs,cal=euler]{mathalfa}
\usepackage{xfrac}
\usepackage[unicode,hidelinks,bookmarks=false]{hyperref}
\newcommand{\ie}{i.e.\ }
\newcommand{\cf}{cf.\ }
\newcommand{\resp}{respectively\ }
\newcommand{\Subsection}{\subsection}
\providecommand{\nobreakdash}{\nobreak\hbox{-}}
\setlength{\emergencystretch}{2em}
\multlinegap0pt
\usepackage{amssymb}
\usepackage{mathtools}
\usepackage{xcolor}
\usepackage{tikz}
\usepackage{algorithm}\usepackage{algpseudocode}
\newtheorem{assumption}{Assumption}
\newtheorem{proposition}{Proposition}
\newtheorem{theorem}{Theorem}
\newtheorem{corollary}{Corollary}
\newcommand{\E}{\mathbb E}
\newcommand{\Pp}{\mathbb P}
\newcommand{\R}{\mathbb R}
\newcommand{\Reach}{\operatorname{Reach}}
\newcommand{\Vol}{\operatorname{Vol}}
\newcommand{\cE}{\mathcal E}
\newcommand{\eps}{\varepsilon}
\renewcommand{\epsilon}{\eps}
\newcommand{\osc}{\operatorname{osc}}
\newcommand{\whatnu}{\widehat\nu_{N,\kappa}}
\newcommand{\wideM}{\widehat M_N}
\title{Low Dimensional Sampling under Reconstructed Constraints}
\author{Imon Banerjee, Riddhiman Bhattacharyya}
\date{Source: arXiv:2609.14837v1 (CC BY 4.0)}
\begin{document}
\maketitle

\noindent\emph{Source and licence.} Low Dimensional Sampling under Reconstructed Constraints. Version arXiv:2609.14837v1 (CC BY 4.0), distributed by arXiv under the Creative Commons Attribution 4.0 International licence, http://creativecommons.org/licenses/by/4.0/, as recorded in the arXiv licence metadata for this exact version and shown on the abstract page https://arxiv.org/abs/2609.14837v1. Full licence notice: \texttt{licenses/arxiv-2609.14837-CC-BY-4.0.txt}.\par\smallskip

\noindent\textit{Editorial scope.} Counted unit: the article's corollary cor:optimized (source0.tex lines 1277--1287). The article's complete original proof of it is delivered with it (source0.tex lines 1710--1750, one \texttt{proof} environment of 41 lines, which the article places away from the statement). No mathematics is written, repaired or re-derived: the statement and the proof are the article's own lines, in the article's own notation, and no retained line is substituted or re-flowed.  Retained uncounted as the source-local context the proof needs: lines 655--668; lines 736--757; lines 760--767; lines 827--844; lines 986--993; lines 997--1025; lines 1055--1097; lines 1109--1131.  Nothing was omitted from any retained span, so the package carries no omission record.  No retained line contains a citation, so the wrapper carries no bibliography and no bibliography entry.  No retained line contains a figure environment, an image-inclusion command or drawing code, so no image is delivered and the package declares no asset dependency.  Editorial formatting only, in addition: an AMS \texttt{amsart} preamble that restates the article's own declarations and macros used by the retained lines, namely the environments \texttt{assumption}, \texttt{proposition}, \texttt{theorem}, \texttt{corollary} on separate counters, together with the standard packages those lines need.  The retained environments print in this document as Assumption 1, Proposition 1, Theorem 1, Corollary 1 in order of appearance, whereas the article prints the same items with the numbering of its own full document.  The complete native source of the article as distributed in the arXiv e-print for this exact version is bundled unchanged as \texttt{sources/210365/source0.tex}.
\begin{assumption}[True manifold and observations]
\label{ass:true-manifold}
For $1\le d<D$, the set $M\subset\R^D$ is a compact, connected,
$d$-dimensional $C^2$ embedded submanifold without boundary, with
\[
  \Reach(M)\ge\tau_0>0.
\]
We assume throughout that $d$ and $\tau_0$ are known. We observe
\[
  X_1,\ldots,X_N\overset{\mathrm{iid}}{\sim}
  \operatorname{Unif}(M)
  :=\frac{\Vol_M}{\Vol_M(M)}.
\]
\end{assumption}

\begin{proposition}[Reconstruction event]
\label{prop:reconstruction-event}
Suppose Assumption~\ref{ass:true-manifold} holds, $1\le d<D$, and the
fixed tuning parameter satisfies $0<\lambda<3^{-1/d}$. There exist
constants $C_H,C_P>0$ and $N_0$ such that, for every $N\ge N_0$,
\begin{equation}
 \Pp\left\{
 d_H(\wideM,M)>
 C_H\left(\frac{\log N}{N}\right)^{2/d}
 \right\}
 \le C_P\frac{(\log N)^{\beta_d}}{N^2}.
\label{eq:divol-reconstruction-tail}
\end{equation}
The constants and threshold may depend on the fixed model, including its
geometry and volume, and on $\lambda$. Consequently, whenever
\eqref{eq:reconstruction-sample-size} holds,
\begin{equation}
  \Pp(\cE_N)\ge1-\alpha_N.
\label{eq:event-probability}
\end{equation}
The probability is over the original observations only.
\end{proposition}

\begin{assumption}[Known ambient enclosure]
\label{ass:ambient-envelope}
There exist known deterministic quantities $z_0\in\R^D$ and $R_0>0$
such that
\[
    M\subseteq \overline B(z_0,R_0).
\]
\end{assumption}

\begin{assumption}[Potential and target distribution]
\label{ass:potential-target}
The potential $g$ is $C^2$ on a neighborhood of $K$. Define the finite
constants and assume they are bounded
\begin{align}
 G&:=\sup_{x\in K}\lVert\nabla g(x)\rVert,\label{eq:G-def}\\
 H&:=\sup_{x\in K}\lVert\nabla^2g(x)\rVert_{\mathrm{op}},\label{eq:H-def}\\
 \osc_K(g)&:=\sup_Kg-\inf_Kg.\label{eq:osc-def}
\end{align}
The desired distribution is the singular Gibbs law
\begin{equation}
  \pi_M(dy)
  :=\frac{e^{-g(y)}}{Z_M}\,\Vol_M(dy),
  \qquad
  Z_M:=\int_Me^{-g(y)}\,\Vol_M(dy).
\label{eq:true-target}
\end{equation}
\end{assumption}

\begin{equation}
 N\ge N_0,
 \qquad
 C_H\left(\frac{\log N}{N}\right)^{2/d}\le\Delta_N,
 \qquad
 C_P\frac{(\log N)^{\beta_d}}{N^2}\le\alpha_N.
\label{eq:reconstruction-sample-size}
\end{equation}

\begin{theorem}[Finite-time RWMH approximation of the manifold target]
\label{thm:rwmh-end-to-end}
Suppose Assumptions~\ref{ass:true-manifold},
\ref{ass:ambient-envelope}, and~\ref{ass:potential-target} hold, and let
$K$ and $r$ be chosen as in Assumption~\ref{ass:ambient-envelope}. Fix $\kappa\ge\kappa_0$ and $\sigma>0$,
and suppose
\[
 0<\Delta_N\le r/4,
 \qquad
 \Delta_N\sqrt\kappa\le1.
\]
Then $\whatnu$ is the unique invariant probability measure of
$P_{N,\kappa}$. There is a fixed-model constant $C_M>0$ such that,
for every initial law $\rho$ on $K$ and every integer $m\ge0$,
on the reconstruction event $\cE_N$,
\begin{align}
 W_2(\rho P_{N,\kappa}^m,\pi_M)
 \le{}&L_K(1-\underline\varepsilon_{\kappa,\sigma})^{m/2}
 \notag\\
 &+C_M\left(
 \kappa^{-1/2}+\Delta_N\sqrt\kappa
 +\kappa^{q/4}e^{-\kappa r^2/64}
 \right).
\label{eq:rwmh-end-to-end}
\end{align}
If the reconstruction sample-size condition
\eqref{eq:reconstruction-sample-size} holds, then this bound holds with
probability at least $1-\alpha_N$.
\end{theorem}

\begin{corollary}[Complexity With Respect to $N$]
\label{cor:rwmh-expected-rate}
Under the hypotheses of Theorem~\ref{thm:rwmh-end-to-end}, let
$0<\epsilon_{\mathrm{tar}}<\min\{1,L_K\}$ be a target Wasserstein accuracy. We use the estimator above and make the analysis choices
\[
 \Delta_N:=\epsilon_{\mathrm{tar}}^2,
 \qquad
 \kappa_N:=\epsilon_{\mathrm{tar}}^{-2},
 \qquad
 \alpha_N:=\epsilon_{\mathrm{tar}}.
\]
Fix $\sigma>0$, and suppose $\epsilon_{\mathrm{tar}}$ is small enough that
$\kappa_N\ge\kappa_0$ and
$\epsilon_{\mathrm{tar}}^2\le r/4$. Suppose $N$ is large enough that the
reconstruction sample-size condition \eqref{eq:reconstruction-sample-size}
holds,
and choose an integer $m_N\ge0$ such that
\begin{equation}
 m_N
 \ge
 \left\lceil
 \frac{2}{\underline\varepsilon_{\kappa_N,\sigma}}
 \log\frac{L_K}{\epsilon_{\mathrm{tar}}}
\right\rceil,
\label{eq:rwmh-optimized-iterations}
\end{equation}
Then, for every initial law $\rho$ on $K$,
\begin{equation}
 \E\left[
 W_2(\rho P_{N,\kappa_N}^{m_N},\pi_M)
 \right]
 =O(\epsilon_{\mathrm{tar}}).
\label{eq:rwmh-expected-rate}
\end{equation}
The implicit constant is independent of $N$ and
$\epsilon_{\mathrm{tar}}$. In particular, suppressing log-factors
\[
 N=\widetilde O\left(
 \epsilon_{\mathrm{tar}}^{-d}
 \right),
\]
is a sufficient observation count for fixed model parameters.
\end{corollary}

\begin{theorem}
\label{thm:main-w2}
Suppose Assumptions~\ref{ass:true-manifold},
\ref{ass:ambient-envelope}, and~\ref{ass:potential-target} hold, and $K$
and $r$ are chosen as above. Suppose
\[
 \kappa\ge\kappa_0,
 \qquad 0<\Delta_N\le r/4,
 \qquad \Delta_N\sqrt\kappa\le1.
\]
There is a fixed-model constant $C_M>0$, independent of
$N,\kappa,\Delta_N$, such that, on $\cE_N$,
\begin{equation}
 W_2(\whatnu,\pi_M)
 \le C_M\left(
 \kappa^{-1/2}+\Delta_N\sqrt\kappa
 +\kappa^{q/4}e^{-\kappa r^2/64}
 \right).
\label{eq:main-w2-bound}
\end{equation}
Consequently, if \eqref{eq:reconstruction-sample-size} holds, this bound
holds with probability at least $1-\alpha_N$ over the observations.
\end{theorem}

\begin{corollary}[Small-error rate]
\label{cor:optimized}
Under the assumptions of Theorem~\ref{thm:main-w2}, choose
$\kappa=\Delta_N^{-1}$ and suppose
$0<\Delta_N\le\min\{1,r/4,\kappa_0^{-1}\}$.
Then, on $\cE_N$,
\begin{equation}
 W_2(\whatnu,\pi_M)=O(\Delta_N^{1/2}).
\label{eq:Delta-half}
\end{equation}
\end{corollary}

\begin{proof}
The inequality $1-u\le e^{-u}$ and
\eqref{eq:rwmh-optimized-iterations} give
\[
 L_K(1-\underline\varepsilon_{\kappa_N,\sigma})^{m_N/2}
 \le\epsilon_{\mathrm{tar}}.
\]
The analysis tolerance is $\Delta_N=\epsilon_{\mathrm{tar}}^2$, and hence
$\kappa_N=\epsilon_{\mathrm{tar}}^{-2}=\Delta_N^{-1}$.
Corollary~\ref{cor:optimized} therefore gives
\[
 W_2(\whatnu,\pi_M)=O(\epsilon_{\mathrm{tar}})
 \qquad\text{on }\cE_N.
\]
The triangle inequality therefore yields
\[
 W_2(\rho P_{N,\kappa_N}^{m_N},\pi_M)
 =O(\epsilon_{\mathrm{tar}})
 \qquad\text{on }\cE_N.
\]
For every reconstruction outcome, both measures in this Wasserstein
distance are supported on $K$, so the distance is at most $L_K$.
For $N$ as in the hypothesis of Corollary~\ref{cor:rwmh-expected-rate},
Proposition~\ref{prop:reconstruction-event} gives
$\Pp(\cE_N^c)\le\epsilon_{\mathrm{tar}}$, and consequently
\[
 \E\left[
 W_2(\rho P_{N,\kappa_N}^{m_N},\pi_M)
 \right]
 \le C\epsilon_{\mathrm{tar}}
   +L_K\Pp(\cE_N^c)
 \le(C+L_K)\epsilon_{\mathrm{tar}}
 =O(\epsilon_{\mathrm{tar}}).
\]
Here $L_K=2(R_0+3r)$ is fixed. For fixed model parameters and sufficiently
small $\epsilon_{\mathrm{tar}}$, a sufficiently large multiple of
$\epsilon_{\mathrm{tar}}^{-d}\log(1/\epsilon_{\mathrm{tar}})$ satisfies
the sample-size hypothesis of Corollary~\ref{cor:rwmh-expected-rate}; the
failure-tail condition is then of smaller order. This proves the stated
sufficient observation count $N=\widetilde O(\epsilon_{\mathrm{tar}}^{-d})$.
\end{proof}

\end{document}
